\setcounter{section}{2}


\inputaufgabe
{}
{





Zeige, dass ein
\definitionsverweis {reelles Geradenbündel}{}{}
\maabb {} { L } { X
} {}
über einem
\definitionsverweis {topologischen Raum}{}{}
$X$ genau dann trivial ist, wenn es einen stetigen nullstellenfreien
\definitionsverweis {Schnitt}{}{}
besitzt.





}
{} {}

\inputaufgabe
{}
{





Es sei
\maabb {s} { X } { V
} {}
ein
\definitionsverweis {stetiger Schnitt}{}{}
zu einem
\definitionsverweis {reellen Vektorbündel}{}{}
\maabb {p} { V } { X
} {}
über einem
\definitionsverweis {topologischen Raum}{}{}
$X$. Zeige, dass das
\definitionsverweis {Bild}{}{}

\mavergleichskette
{\vergleichskette
{s(X)
}
{ \subseteq }{ V
}
{  }{
}
{    }{
}
{   }{
}
}
{}{}{}
eine
\definitionsverweis {abgeschlossene Teilmenge}{}{}
ist, die
\definitionsverweis {homöomorph}{}{}
zu $X$ ist.





}
{} {}

\inputaufgabe
{}
{





Es sei
\maabb {p} { V } { U
} {}
ein Vektorbündel über

\mavergleichskette
{\vergleichskette
{U
}
{ \subseteq }{\R^m
}
{  }{
}
{    }{
}
{   }{
}
}
{}{}{,}
offen, das durch ein
\definitionsverweis {lineares Gleichungssystem}{}{}
in $n$ Variablen mit $m$ universellen Parametern gegeben ist, wobei $U$ dadurch gekennzeichnet sei, dass die Dimension der Fasern konstant sei. Zeige, dass ein
\definitionsverweis {stetiger Schnitt}{}{}
in $V$ das gleiche ist wie eine stetige universelle Lösung des linearen Gleichungssystems.





}
{} {}

\inputaufgabegibtloesung
{}
{





Man gebe ein
\definitionsverweis {stetiges}{}{} \definitionsverweis {Vektorfeld}{}{}
auf $S^2$ an, das nur eine Nullstelle besitzt.





}
{} {}

\inputaufgabe
{}
{





Zeige, dass die Einschränkung des Vektorbündels aus

Beispiel 1.2

auf die offene Teilmenge

\mavergleichskettedisp
{\vergleichskette
{D(r) \cup D(s)
}
{ =} { \left\{ (r,s,t)  \mid  r \neq 0 \text{ oder } s \neq 0        \right\}
}
{ \subset} {\R^3
}
{ } {
}
{ } {
}
}
{}{}{}
eine Trivialisierung besitzt.





}
{} {}

\inputaufgabe
{}
{





Es sei

\mavergleichskette
{\vergleichskette
{X
}
{ = }{ \bigcup_{i \in I} U_i
}
{  }{
}
{    }{
}
{   }{
}
}
{}{}{}
eine
\definitionsverweis {offene Überdeckung}{}{}
eines
\definitionsverweis {topologischen Raumes}{}{}
$X$. Inwiefern wird dadurch ein
\definitionsverweis {topologisches Verklebungsdatum}{}{}
gegeben? Lässt sich aus dem Verklebungsdatum der topologische Raum $X$ rekonstruieren?





}
{} {}

\inputaufgabe
{}
{





Man mache sich klar, wodurch ein Verklebungsdatum mit zwei offenen Mengen gegeben ist.





}
{} {}

\inputaufgabe
{}
{





Man mache sich klar, wodurch ein Verklebungsdatum mit drei offenen Mengen gegeben ist.





}
{} {}

\inputaufgabe
{}
{





Es sei
\mathkor {} {U} {und} {V} {}
jeweils eine reelle Gerade, und diese werden entlang der offenen Halbgeraden

\mavergleichskette
{\vergleichskette
{\R_+
}
{ \subseteq }{ U
}
{  }{
}
{    }{
}
{   }{
}
}
{}{}{}
und

\mavergleichskette
{\vergleichskette
{\R_+
}
{ \subseteq }{ V
}
{  }{
}
{    }{
}
{   }{
}
}
{}{}{}
verklebt. Ist der entstehende Raum
\definitionsverweis {Hausdorffsch}{}{?}





}
{} {}

\inputaufgabe
{}
{





Es sei
\mathkor {} {U} {und} {V} {}
jeweils eine reelle Gerade, und diese werden entlang der punktierten Geraden

\mavergleichskette
{\vergleichskette
{\R \setminus \{0\}
}
{ \subseteq }{ U
}
{  }{
}
{    }{
}
{   }{
}
}
{}{}{}
und

\mavergleichskette
{\vergleichskette
{\R \setminus \{0\}
}
{ \subseteq }{ V
}
{  }{
}
{    }{
}
{   }{
}
}
{}{}{}
miteinander verklebt. Ist der entstehende Raum
\definitionsverweis {Hausdorffsch}{}{?}





}
{} {}

\inputaufgabe
{}
{





Wir betrachten den
\definitionsverweis {topologischen Raum}{}{}

\mavergleichskettedisp
{\vergleichskette
{X
}
{ =} { \R \cup \{ 0' \}
}
{ } {
}
{ } {
}
{ } {
}
}
{}{}{,}
der aus den reellen Zahlen entsteht, indem man ein neues Element $0'$  hinzunimmt, und die folgenden zwei Arten von Teilmengen für offen erklärt: Die Mengen

\mavergleichskette
{\vergleichskette
{U
}
{ \subseteq }{ \R
}
{  }{
}
{    }{
}
{   }{
}
}
{}{}{}
offen mit

\mavergleichskette
{\vergleichskette
{ 0
}
{ \notin }{U
}
{  }{
}
{    }{
}
{   }{
}
}
{}{}{}
und die Mengen
\mathl{V \cup \{0'\}}{} für

\mavergleichskette
{\vergleichskette
{V
}
{ \subseteq }{\R
}
{  }{
}
{    }{
}
{   }{
}
}
{}{}{}
offen mit

\mavergleichskette
{\vergleichskette
{ 0
}
{ \in }{ V
}
{  }{
}
{    }{
}
{   }{
}
}
{}{}{.}
Zeige, dass $X$ ein topologischer Raum ist. Ist der Punkt $0$ in diesem Raum abgeschlossen? Ist der Punkt $0'$ in diesem Raum abgeschlossen? Wie verhält sich dieser Raum zu dem in

Aufgabe 2.10

konstruierten Raum?





}
{} {}

\inputaufgabe
{}
{





Es sei
\mathkor {} {U} {und} {V} {}
jeweils eine reelle Gerade, und diese werden entlang der punktierten Geraden

\mavergleichskette
{\vergleichskette
{\R \setminus \{0\}
}
{ \subseteq }{ U
}
{  }{
}
{    }{
}
{   }{
}
}
{}{}{}
und

\mavergleichskette
{\vergleichskette
{\R \setminus \{0\}
}
{ \subseteq }{ V
}
{  }{
}
{    }{
}
{   }{
}
}
{}{}{}
mithilfe der inversen Abbildung
\mathl{t \mapsto t^{-1}}{} verklebt. Welcher topologische Raum entsteht dabei?





}
{} {}

\inputaufgabe
{}
{





Es sei
\mathkor {} {U} {und} {V} {}
jeweils eine komplexe Gerade ${\mathbb C}$
\zusatzklammer {also die Gaußsche Zahlenebene} {} {}
und diese werden entlang der gelochten Ebene

\mavergleichskette
{\vergleichskette
{ {\mathbb C}  \setminus \{0\}
}
{ \subseteq }{ U
}
{  }{
}
{    }{
}
{   }{
}
}
{}{}{}
und

\mavergleichskette
{\vergleichskette
{ {\mathbb C}  \setminus \{0\}
}
{ \subseteq }{ V
}
{  }{
}
{    }{
}
{   }{
}
}
{}{}{}
mithilfe der inversen Abbildung
\mathl{z \mapsto z^{-1}}{} verklebt. Welcher topologische Raum entsteht dabei?





}
{} {}

\inputaufgabe
{}
{





Zeige, dass im Beweis zu

Lemma 2.6

in der Tat eine
\definitionsverweis {Äquivalenzrelation}{}{}
vorliegt.





}
{} {}

\inputaufgabe
{}
{





Es sei ein
\definitionsverweis {Verklebungsdatum}{}{}
\mathkor {} {U_i} {} {i \in I} {,}
mit

\mavergleichskette
{\vergleichskette
{U_{ij}
}
{ = }{ \emptyset
}
{  }{
}
{    }{
}
{   }{
}
}
{}{}{}
für alle

\mavergleichskette
{\vergleichskette
{i
}
{ \neq }{j
}
{  }{
}
{    }{
}
{   }{
}
}
{}{}{}
gegeben. Welchen topologischen Raum legt dieses Verklebungsdatum fest?





}
{} {}

\inputaufgabe
{}
{





Wir betrachten den Zylinder

\mavergleichskettedisp
{\vergleichskette
{U
}
{ =} { V
}
{ =} { S^1 \times {]0,3[}
}
{ } {
}
{ } {
}
}
{}{}{.}
Auf den offenen Teilmengen
\mathl{S^1 \times {]0,1[}}{} bzw.
\mathl{S^1 \times {]2,3[}}{}
\zusatzklammer {in $U$ bzw. in $V$} {} {}
betrachten wir die
\definitionsverweis {Homöomorphismen}{}{,}
die sich aus der Identität auf dem Kreis bzw. der Punktspiegelung einerseits und der Identität auf dem Intervall bzw. der Spiegelung in der Intervallmitte andererseits ergibt. Welche geometrischen Objekte ergeben sich durch diese verschiedenen
\definitionsverweis {Verklebungsdaten}{}{?}





}
{} {}

\inputaufgabe
{}
{





Es sei $M$ eine
\definitionsverweis {topologische Mannigfaltigkeit}{}{}
und
\maabbdisp {\alpha_i} {U_i } { V_i
} {}
eine Familie von
\definitionsverweis {Karten}{}{}
mit den
\definitionsverweis {Übergangsabbildungen}{}{}
\maabbdisp {\varphi_{ij} = \alpha_j \circ (\alpha_i)^{-1}} {V_i \cap \alpha_i(U_i \cap U_j) } { V_j \cap \alpha_j(U_i \cap U_j)
} {.}
Zeige, dass man aus der Familie der

\mathbed {V_i} {}
 {i \in I} {}
 {} {} {} {,}
den Teilmengen
\mathl{V_{ij} \subseteq V_i}{} und den Übergangsabbildungen
\maabbdisp {\varphi_{ij}} {V_{ij}} { V_{ji}
} {}
die Mannigfaltigkeit $M$ rekonstruieren kann.

a) Betrachte auf

\mavergleichskettedisp
{\vergleichskette
{N
}
{ \defeq} { \biguplus_{i \in I} V_i
}
{ } {
}
{ } {
}
{ } {
}
}
{}{}{}
die
\definitionsverweis {Äquivalenzrelation}{}{,}
unter der zwei Punkte
\mathl{P\in V_i}{} und
\mathl{Q \in V_j}{} gleich sind, wenn sie unter
\mathl{\varphi_{ij}}{} ineinander abgebildet werden.

b) Versehe die Quotientenmenge
\mathl{N/ \sim}{} mit einer geeigneten Topologie.

c) Definiere auf
\mathl{N/ \sim}{} Karten.

d) Zeige, dass
\mathkor {} {M} {und} {N/ \sim} {}
\definitionsverweis {homöomorph}{}{}
sind.





}
{} {}

\inputaufgabe
{}
{





Es sei ein
\definitionsverweis {Verklebungsdatum}{}{}

\mathbed {U_i} {}
 {i \in I} {}
 {} {} {} {,}
für
\definitionsverweis {topologische Räume}{}{}
gegeben. Es sei $Z$ ein weiterer topologischer Raum und es seien
\definitionsverweis {stetige Abbildungen}{}{}

\maabbdisp {\theta_i} {U_i } { Z
} {}
gegeben, die die Bedingung

\mavergleichskette
{\vergleichskette
{ \theta_i \vert_{U_{ij} }
}
{ = }{ \left(  \theta_j \vert_{U_{ji} }  \right)  \circ \varphi_{ji}
}
{  }{
}
{    }{
}
{   }{
}
}
{}{}{}
erfüllen. Zeige, dass es dann eine eindeutig bestimmte stetige Abbildung
\maabbdisp {\theta} { X } { Z
} {}
mit

\mavergleichskette
{\vergleichskette
{ \left(  \psi_i  \right)^{-1} \circ \theta \vert_{V_i }
}
{ = }{ \theta_i
}
{  }{
}
{    }{
}
{   }{
}
}
{}{}{}
gibt, wobei $X$ den durch die Verklebungsdaten festgelegten topologischen Raum
\zusatzklammer {siehe

Lemma 2.6
,
auch für die Notation} {} {}
bezeichnet.





}
{} {}

\inputaufgabe
{}
{





Zeige, dass ein Verklebungsdatum für ein Geradenbündel zu einer offenen Überdeckung

\mavergleichskette
{\vergleichskette
{X
}
{ = }{ \bigcup_{i \in I} U_i
}
{  }{
}
{    }{
}
{   }{
}
}
{}{}{}
das gleiche ist wie eine Familie von stetigen nullstellenfreien Funktionen
\maabb {f_{ij}} {U_i \cap U_j } { \R
} {,}
die auf
\mathl{U_i \cap U_j \cap U_k}{} jeweils die Bedingung

\mavergleichskette
{\vergleichskette
{ f_{ki}
}
{ = }{ f_{kj} \cdot f_{ji}
}
{  }{
}
{    }{
}
{   }{
}
}
{}{}{}
erfüllen.





}
{} {}

\inputaufgabe
{}
{





Bestimme
\definitionsverweis {Verklebungsdaten}{}{}
für das Vektorbündel aus

Beispiel 1.2
.





}
{} {}

\inputaufgabe
{}
{





Es sei

\mavergleichskette
{\vergleichskette
{X
}
{ = }{ \bigcup_{i \in I} U_i
}
{  }{
}
{    }{
}
{   }{
}
}
{}{}{}
eine
\definitionsverweis {offene Überdeckung}{}{}
eines
\definitionsverweis {topologischen Raumes}{}{}
$X$. Es seien
\maabbdisp {\varphi_{ij}} {U_i \cap U_j } { \operatorname{GL}_{  r   } \! \left(  \R  \right)
} {}
und
\maabbdisp {\psi_{ij}} {U_i \cap U_j } { \operatorname{GL}_{  r   } \! \left(  \R  \right)
} {}
\definitionsverweis {Matrixbeschreibungen}{}{,}
die zu den Vektorbündeln
\mathkor {} {E} {bzw.} {F} {}
führen. Zeige, dass diese Bündel genau dann
\definitionsverweis {isomorph}{}{}
sind, wenn es stetige Abbildungen
\maabbdisp {\alpha_i} { U_i } { \operatorname{GL}_{  r   } \! \left(  \R  \right)
} {}
derart gibt, dass
\zusatzklammer {für sinnvolle Einschränkungen} {} {}

\mavergleichskettedisp
{\vergleichskette
{ \alpha_i \circ  \varphi_{ij}  \circ \alpha_j^{-1}
}
{ =} { \psi_{ij}
}
{ } {
}
{ } {
}
{ } {
}
}
{}{}{}
für alle $i,j$ gilt.





}
{} {}

\inputaufgabe
{}
{





Ein
\definitionsverweis {Vektorbündel}{}{}
\maabb {} { V } { X
} {}
vom Rang $m$ über einem
\definitionsverweis {topologischen Raum}{}{}
$X$ sei durch stetige Matrizenabbildungen
\maabbdisp {\varphi_{ij}} { U_i \cap U_j } { \operatorname{GL}_{  m   } \! \left(  \R  \right)
} {}
zu einer offenen Überdeckung

\mavergleichskette
{\vergleichskette
{X
}
{ = }{ \bigcup_{i \in I} U_i
}
{  }{
}
{    }{
}
{   }{
}
}
{}{}{}
gegeben. Zeige, dass ein
\definitionsverweis {stetiger Schnitt}{}{}
\maabbdisp {s} { X } { V
} {}
dasselbe ist wie eine Familie von stetigen Abbildungen
\maabb {t_i} {U_i } { \R^m
} {,}
die die Bedingungen

\mavergleichskettedisp
{\vergleichskette
{t_j
}
{ =} { \varphi_{ij} t_i
}
{ } {
}
{ } {
}
{ } {
}
}
{}{}{}
für alle $ij$ erfüllt.





}
{} {}

\inputaufgabe
{}
{





Wir betrachten die algebraische Realisierung des Möbiusbandes aus

Beispiel 2.12
,
also

\mavergleichskettedisp
{\vergleichskette
{Y
}
{ =} { \left\{ (x,y,z,w) \in \R^4  \mid   x^2+y^2 = 1 , \, (1-y)z = xw   , \,  xz = (1+y)w    \right\}
}
{ \rightarrow} { S^1
}
{ } {
}
{ } {
}
}
{}{}{.}
Zeige, dass das Bild der stetigen Abbildung
\maabbeledisp {\psi} { [0, 2\pi] } { \R^4
} { t } { \left(  \cos  t   , \,   \sin  t  , \,   \cos  { \frac{   t  }{  2 } } , \,   \sin  { \frac{   t  }{  2 } }                \right)
} {,}
in $Y$ landet, dass
\mathl{p \circ \psi}{} die trigonometrische Parametrisierung des Einheitskreises ist und dass das Bild von $\psi$ niemals den Nullschnitt trifft.





}
{} {}

Die folgende Aussage lässt sich für die abgeschlossene Version des Möbiusbandes einfach mit der Schere bestätigen.

\inputaufgabe
{}
{





Folgere aus

Aufgabe 2.23
,
dass das Komplement des Nullschnittes im Möbiusband
\definitionsverweis {wegzusammenhängend}{}{}
ist.





}
{} {}

\inputaufgabe
{}
{





Zeige, dass das Komplement des Nullschnittes in einem trivialen
\definitionsverweis {Geradenbündel}{}{}
nicht
\definitionsverweis {zusammenhängend}{}{}
ist.





}
{} {}

\inputaufgabe
{}
{





Man nehme ein schmales rechteckiges Band, verdrehe es mit einer Volldrehung um die längere Achse und verklebe die beiden kürzeren Ränder. Nun schneide man mit einer Schere das Band längs in  der Mitte durch. Ist das entstehende Objekt
\definitionsverweis {zusammenhängend}{}{?}
Wie sieht es aus, wenn man $n$ Halbdrehungen macht?





}
{} {}

\inputaufgabe
{}
{





Bestimme den
\definitionsverweis {Grenzwert}{}{}
der Funktion

\mavergleichskette
{\vergleichskette
{ { \frac{   1-y }{  x } }
}
{ = }{ { \frac{   x  }{  1+y } }
}
{  }{
}
{    }{
}
{   }{
}
}
{}{}{}
auf dem
\definitionsverweis {Einheitskreis}{}{}
für
\mathl{(x,y) \rightarrow (0,1)}{.} Man betrachte dazu auch die trigonometrische Parametrisierung des Einheitskreises.





}
{} {}
